% =====================================================================
%  Block-selection strategies for hourly market probing:
%  mathematical principles and code (development note, pre-lock)
%
%  GENERATED FILE: do not edit; edit strategies_template.tex, then run
%      python build_strategies_tex.py
%  The code listings are copied from the study's source files at build
%  time, with their real line numbers.
%  Compile with pdfLaTeX (e.g. on Overleaf); standard packages only.
% =====================================================================
\documentclass[11pt]{article}
\usepackage[a4paper,margin=2.2cm]{geometry}
\usepackage{amsmath,amssymb,amsthm}
\usepackage{bm}
\usepackage{booktabs}
\usepackage{array}
\usepackage{tabularx}
\usepackage{xcolor}
\usepackage{listings}
\usepackage[hidelinks]{hyperref}

\newtheorem{lemma}{Lemma}
\newtheorem{proposition}{Proposition}
\theoremstyle{definition}
\newtheorem{definition}{Definition}
\newtheorem{remark}{Remark}

\newcommand{\R}{\mathbb{R}}
\newcommand{\tp}{^{\top}}
\newcommand{\one}{\mathbf{1}}
\newcommand{\bu}{\bm{u}}
\newcommand{\bs}{\bm{s}}
\newcommand{\cD}{\mathcal{D}}
\newcommand{\cU}{\mathcal{U}}
\newcommand{\cP}{\mathcal{P}}
\newcommand{\cQ}{\mathcal{Q}}
\newcommand{\code}[1]{\texttt{#1}}

\lstset{
  language=Python,
  basicstyle=\ttfamily\scriptsize,
  keywordstyle=\color{blue!60!black},
  commentstyle=\color{green!40!black},
  stringstyle=\color{red!50!black},
  numbers=left,
  numberstyle=\tiny\color{gray},
  frame=single,
  breaklines=true,
  breakatwhitespace=false,
  columns=fullflexible,
  keepspaces=true,
  showstringspaces=false,
  tabsize=4,
  upquote=true
}

\begin{document}

\title{Block-Selection Strategies for Hourly Market Probing:\\ Mathematical Principles and Code}
\author{Development note of the hourly probing study (pre-lock, not evidence)}
\date{2026-11-23}
\maketitle

\begin{abstract}
At every market hour the operator may buy one block of active power from household batteries and EV chargers, observe the metered voltages, and refit every candidate topology. This note states, in one notation, how each strategy tested in the development runs chooses that block: passive, random, the crude greedy, query by committee, the v3 dictionary planner, the certified planner and the strong greedy. All strategies share the scoring, stopping and safety machinery of Section~\ref{sec:common}; they differ only in the block they buy. Section~\ref{sec:V} explains the price of certainty $V$ and checks the strong greedy's risk bound against the logged runs. Appendix~\ref{app:code} lists the code that implements every formula, copied from the source files with their line numbers.
\end{abstract}

\tableofcontents

% =====================================================================
\section{Setting and notation}\label{sec:setting}

\begin{itemize}
  \item Metered terminals $i=1,\dots,n$ (root excluded); 15-minute slots $t$; market hour $h$ holds slots $4h,\dots,4h+3$.
  \item Measured per slot: net consumption $p_t,q_t\in\R^{n}$ (pu), the root's squared voltage $v_{0,t}^2$, and the squared-voltage drop
        $y_{t,i}=v_{0,t}^2-v_{t,i}^2$.
  \item \textbf{Candidate topology} $k$: a laminar family $\mathcal S_k$ of supports $C\subseteq\{1,\dots,n\}$ (the terminals downstream of a branch), with incidence $Z_k\in\{0,1\}^{S_k\times n}$, $(Z_k)_{C,i}=1$ if $i\in C$. Support flows $a_{t,C}=\sum_{j\in C}p_{t,j}$ and $b_{t,C}=\sum_{j\in C}q_{t,j}$.
  \item \textbf{Response model} (branch-flow, loss-compensated), parameters $\theta_k=(r_C,x_C)_{C\in\mathcal S_k}$ with $r_C=2R_C/Z_{\rm base}$:
  \begin{equation}\label{eq:model}
    y_{t,i}=\sum_{C\ni i}\big(r_C\,a_{t,C}+x_C\,b_{t,C}\big)+\lambda_{k,t,i}(\theta_k)+o_i+\varepsilon_{t,i},
    \qquad\text{in matrix form}\quad y=\Phi_k\theta_k+\lambda_k(\theta_k)+(\one\otimes o)+\varepsilon ,
  \end{equation}
  where $\Phi_k$ is the design matrix, $\lambda_k$ the loss term and $o\in\R^n$ the per-meter offsets (revised step~1).
  \item \textbf{Block} $\bu\in\R^{n}$ (kW, held constant for one hour) is added to the metered consumption; $\bu>0$ is extra consumption (charging), $\bu<0$ a reduction (discharge or paused charging).
  \item \textbf{Timing.} Gate $g$ is the start of hour $g$. It sees data up to the end of hour $g-1$ and seals the block for hour $g+1$ (gate lead 1~h). The block for hour $g$ is already sealed and still to run.
\end{itemize}

% =====================================================================
\section{Machinery shared by every strategy}\label{sec:common}

\subsection{GLS whitening for the shared root meter}
Every terminal's $y$ subtracts the same root reading, so the noise of one slot has covariance proportional to $I+\rho\,\one\one\tp$ with $\rho=\overline{v_0^4}/\overline{v^4}$. Its inverse square root is
\begin{equation}\label{eq:gls}
  W=I-\gamma\,\tfrac{1}{n}\one\one\tp,\qquad \gamma=1-\frac{1}{\sqrt{1+n\rho}},
\end{equation}
applied to every slot (code: \code{model.root\_gamma}, \code{model.whiten\_slots}, \code{Gate.wh}). Every residual, design and difference below is whitened.

\subsection{Hourly refit}
With $D$ the per-terminal demeaning over time (a free offset per meter, by Frisch--Waugh--Lovell), each candidate is refitted on all data so far by bounded least squares,
\begin{equation}\label{eq:refit}
  \hat\theta_k=\arg\min_{\ell_k\le\theta\le\bar\theta_k}\big\|DW\big(y-\lambda_k(\hat\theta_k^{\rm prev})\big)-DW\Phi_k\,\theta\big\|^2,
  \qquad \mathrm{SSE}_k=\text{its minimum},
\end{equation}
with $r_C\ge r_{\min}$ on non-singleton supports. The loss term uses the previous fit (one fixed-point step per gate, three at the first). The parameter covariance on the free coordinates $F$ (those not at a bound) is
\begin{equation}\label{eq:cov}
  \Sigma_k=\hat\sigma^2\big(A_{k,F}\tp A_{k,F}\big)^{-1},\qquad A_k=DW\Phi_k ,
\end{equation}
and zero on the bound coordinates (code: \code{Fit.refit}, \code{Fit.covariance}).

\subsection{Scores, posterior and operational equivalence}
With $N$ rows of data and $p_k=2S_k+n$ parameters (branch $r,x$ and the $n$ offsets),
\begin{equation}\label{eq:bic}
  J_k=\frac{\mathrm{SSE}_k}{2\hat\sigma^2}+\frac{p_k}{2}\ln N,\qquad \hat\sigma^2=\min_k\frac{\mathrm{SSE}_k}{N-p_k},
  \qquad L=\arg\min_k J_k,\qquad E_k=J_k-J_L\ge0,
\end{equation}
and the posterior weights are $w_k=e^{-E_k}\big/\sum_l e^{-E_l}$.

\begin{definition}[Operational equivalence]
Candidate $k$ with $E_k<25$ is equivalent to the leader if its predicted voltages over the forecast day, $\hat v_{k,t,i}=\sqrt{v_{0,t}^2-\hat y_{k,t,i}}$, satisfy $\max_{t,i}|\hat v_{k,t,i}-\hat v_{L,t,i}|\le\varepsilon_V=5\cdot10^{-4}$~pu. $\cQ_L$ is the leader's equivalence class; the \emph{rivals} $R=\{k\notin\cQ_L\}$ are sorted by $E_k$.
\end{definition}

\subsection{Confirmation, stopping and the trust gate}
\begin{itemize}
  \item \textbf{Confirmation.} With $\tau=\alpha/(1-\alpha)=1/19$ ($\alpha=0.05$), gate $g$ passes when
  \begin{equation}\label{eq:confirm}
    \sum_{j\in R}e^{-E_j}\le\tau ,
  \end{equation}
  that is, the leader's class holds posterior mass of at least 0.95. The leader is \emph{confirmed} when the same leader passes at 2 consecutive gates.
  \item \textbf{When a block may be bought.} The leader is not yet confirmed and \eqref{eq:confirm} fails, $R\neq\emptyset$, the no-progress rule has not stopped the run, and the \emph{trust gate} passes: the bank's closure finished without a cap binding, and the leader's adequacy ratio
  \begin{equation}\label{eq:adequacy}
    a_L=\frac{\mathrm{SSE}_L/(N-p_L)}{c\,\sigma_V^2\,1.02^4}\le 1.2,\qquad c=4\text{ with GLS ($8$ without)},
  \end{equation}
  where $c\,\sigma_V^2\,1.02^4$ is the meter-class variance of $y$. In the fresh-seed test (DEVELOPMENT\_RUNS \S18) $a_L$ had the same distribution with and without the truth in the bank (median 0.97), so the gate barely restricted probing on truth-removed banks (55\% of days against 58\%). A near-twin of the truth fits the meters as well as the truth does.
  \item \textbf{No-progress rule.} At gates where a block was bought, record $G_g=\min_{j\in A}E_j$ over the strategy's active rivals $A$. Stop probing if $G$ has risen by less than 1 over the last 6 such gates.
\end{itemize}

\subsection{Predictions for the block hour and pairwise separation}
For the next hour (4 slots), candidate $k$ predicts $\mu_k\in\R^{4n}$ under the load forecast, with the operating-point sensitivity $M_k=\partial y/\partial p\in\R^{4n\times n}$ ($R_k=Z_k\tp\,\mathrm{diag}(r)Z_k$ plus the loss term's Jacobian). Under a block,
\begin{equation}\label{eq:mu}
  \mu_k(\bu)=\mu_k+M_k\,\bu/P_{\rm base}.
\end{equation}
The \textbf{separation} of two candidates $a,b$ is the whitened Mahalanobis distance of their predictions,
\begin{equation}\label{eq:q}
  q_{ab}(\bu)=\big(W\Delta\mu_{ab}(\bu)\big)\tp\bar S_{ab}^{-1}\big(W\Delta\mu_{ab}(\bu)\big),\quad
  \Delta\mu_{ab}(\bu)=\mu_a(\bu)-\mu_b(\bu),\quad
  \bar S_{ab}=\hat\sigma^2I+\tfrac12\tilde\Phi_a\Sigma_a\tilde\Phi_a\tp+\tfrac12\tilde\Phi_b\Sigma_b\tilde\Phi_b\tp ,
\end{equation}
with $\tilde\Phi=W\Phi(P_f,Q_f)$ at the forecast operating point. The terms $\tilde\Phi\Sigma\tilde\Phi\tp$ count the unknown impedances as noise, so a rival that can re-tune its $r,x$ to mimic the leader gets little separation. For a block $\bu=\kappa\bs$ the separation is a convex quadratic in the level,
\begin{equation}\label{eq:quad}
  q_{ab}(\kappa\bs)=q^{(0)}_{ab}+2q^{(1)}_{ab}(\bs)\,\kappa+q^{(2)}_{ab}(\bs)\,\kappa^2,\qquad
  q^{(0)}=d\tp\bar S^{-1}d,\ \ q^{(1)}=e\tp\bar S^{-1}d,\ \ q^{(2)}=e\tp\bar S^{-1}e\ge0,
\end{equation}
with $d=W\Delta\mu_{ab}(0)$ and $e=W(M_a-M_b)\bs/P_{\rm base}$. One Cholesky factor of $\bar S_{ab}$ per pair serves every shape and level.

\subsection{Admissible blocks: dictionary, ceilings and cost}\label{sec:window}
\textbf{Dictionary.} $\cD$ contains the single-home shapes $\pm\sigma_i e_i$ ($\sigma_i$: standard deviation of terminal $i$'s metered load, kW) and, for each of the three nearest rivals $j$, the group shapes $\pm\sigma\odot\one_{N_j}$ on the terminals $N_j$ whose non-singleton supports differ between $L$ and $j$ (when $1<|N_j|<n$). $\cD_1\subset\cD$ denotes the single-home shapes.

\textbf{Ceiling.} A shape $\bs$ may be bought at level $0<\kappa\le\bar\kappa(\bs)$ with
\begin{equation}\label{eq:ceiling}
  \bar\kappa(\bs)=\min\{\kappa_A,\ \kappa_V,\ \kappa_S,\ \kappa_U,\ \kappa_L,\ \kappa_B\},
\end{equation}
taken over the $H_r=\min(2,24-h)$ planning hours:
\begin{itemize}
  \item $\kappa_A$, \emph{headroom and energy}: $\min_{i:s_i\neq0}\bar u_i^{\pm}/|s_i|$, with
        $\bar u_i^{+}=\min\!\big(P^{\rm bat}_i,(\mathrm{soc}_{\max}-\mathrm{soc}_i)E_i^{\rm bat}/(\eta H_r)\big)+P^{\rm EV}_i\,\one[\text{plugged}]$ and
        $\bar u_i^{-}=\min\!\big(P^{\rm bat}_i,(\mathrm{soc}_i-\mathrm{soc}_{\min})E_i^{\rm bat}\eta/H_r\big)+0.5\,P^{\rm EV}_i\,\one[\text{plugged}]$.
        The state of charge includes the block already sealed for the current hour.
  \item $\kappa_V$, \emph{robust voltage}: over the plausible candidates and every slot and terminal, the voltage under the forecast load scaled by $1\pm2\sigma_f$ ($\sigma_f=8\%$), moved by the block's first-order effect $g\kappa$ plus $z\,b\kappa$ for the impedance uncertainty ($z=2.33$), stays in $[0.96,1.035]$~pu.
  \item $\kappa_S$, \emph{transformer}: $|P_0+\kappa\one\tp\bs|\le\sqrt{S_{\max}^2-Q_0^2}$ in every slot.
  \item $\kappa_U=\kappa_V/6$, \emph{single-phase unbalance}.
  \item $\kappa_L$, \emph{linearisation}: the second-order remainder of the block stays below the meter noise $\sigma$.
  \item $\kappa_B$, \emph{budget}: $c(\kappa\bs)\le\min\big(c_{\rm hour},\,(c_{\rm day}-c_{\rm committed})/H_r\big)$, where $c_{\rm committed}$ is the day's spend including the block already sealed for the current hour.
\end{itemize}

\textbf{Cost.} The development runs since DEVELOPMENT\_RUNS \S16 settle every block as demand response, per kWh moved in the hour:
\begin{equation}\label{eq:cost}
  c(\bu)=\pi_\downarrow\sum_i\max(-u_i,0)+\pi_\uparrow\sum_i\max(u_i,0),\qquad \pi_\downarrow=4,\ \pi_\uparrow=1.2\ \text{CNY/kWh},
\end{equation}
with caps $c_{\rm hour}=20$ and $c_{\rm day}=100$~CNY. The v3 draft contract $c(\bu)=\sum_i\pi_i(h)\,u_i^2/\bar u_i$ is kept as an option. Both are homogeneous in the level, $c(\kappa\bs)=\kappa^{a}c(\bs)$ with $a=1$ (linear) or $a=2$ (contract), so the budget level is $\kappa_B=(c_{\max}/c(\bs))^{1/a}$.

\textbf{Level search.} Every planner that optimises a level scans the 24 levels $\kappa\in\{\bar\kappa/24,\,2\bar\kappa/24,\dots,\bar\kappa\}$ for every admissible shape.

% =====================================================================
\section{The strategies}\label{sec:strategies}

Every strategy below runs only at gates where a block may be bought (Section~\ref{sec:common}); otherwise it buys nothing. ``Greedy'' means deciding one hour at a time: at each gate, buy the block with the best value for the next hour given everything measured so far, without planning later hours.

\subsection{Passive (\code{passive\_bic})}
$\bu_g\equiv0$. The refit, scores and confirmation rule still run every hour.

\subsection{Random, matched (\code{random\_matched\_P})}
Replays the hours $h$ and the hourly costs $\bar c_h$ of the certified planner on the same input and path. At each such hour it draws a terminal $i$ and a sign $\varsigma\in\{+1,-1\}$ uniformly and buys
\begin{equation}
  \bu=\varsigma\,\min\!\big(\bar c_h/\pi_\varsigma,\ \bar u_i^{\varsigma}\big)\,e_i ,
\end{equation}
with $\pi_{+}=\pi_\uparrow$ and $\pi_{-}=\pi_\downarrow$. It is matched on hours and target spend, not targeted. When the headroom binds it spends less than the certified planner. The cap $\bar u_i^{\varsigma}$ is the state-of-charge-aware headroom used by every other arm. Before 2026-11-23 the upward cap ignored a full battery: in DEVELOPMENT\_RUNS \S18, 197 of 532 random runs delivered more than the batteries and EVs could give, most likely to the random arm's advantage. Rerun with the fixed cap, the random arm exceeded the headroom in no run and was correct in 187 runs against 190, and the strong greedy beat it 50 to 8 instead of 47 to 8.

\subsection{Crude greedy (\code{greedy\_P})}
Against the nearest rival $j_1=\arg\min_{j\in R}E_j$, at full ceiling, single homes only:
\begin{equation}
  \bs^\star=\arg\max_{\bs\in\cD_1}\big\|\mu_L(\bar\kappa(\bs)\bs)-\mu_{j_1}(\bar\kappa(\bs)\bs)\big\|^2,\qquad \bu=\bar\kappa(\bs^\star)\,\bs^\star .
\end{equation}
It ignores the other rivals, the shared root error, parameter uncertainty and cost.

\subsection{Query by committee (\code{qbc\_P})}
The committee $\cP$ is the up to 8 most probable candidates with $w_k\ge0.01$ (the leader always included), with renormalised weights $\bar w_k$:
\begin{equation}
  \bs^\star=\arg\max_{\bs\in\cD_1}\sum_{k\in\cP}\bar w_k\big\|\mu_k(\bu_{\bs})-\bar\mu(\bu_{\bs})\big\|^2,\qquad
  \bar\mu(\bu)=\sum_{k\in\cP}\bar w_k\,\mu_k(\bu),\qquad \bu_{\bs}=\bar\kappa(\bs)\,\bs .
\end{equation}
It buys at the ceiling whenever a block may be bought. It has no stopping rule of its own, no whitening, no parameter uncertainty and no notion of operational equivalence, so it also pays to separate a leader from its own twins.

\subsection{v3 dictionary planner (\code{t\_opt\_P}, PROTOCOL v3)}
With Chernoff order $\omega=\tfrac12$, the v3 planner requires, for every rival $j$ in an active set $A$,
\begin{equation}\label{eq:v3}
  q_j(\bu)\ge\eta_j=\max\!\Big(\frac{2\gamma_j}{\omega(1-\omega)},0\Big),\qquad
  \gamma_j=\omega\big(\ln\tfrac{1}{\tau}-E_j\big)+\ln\tfrac{1}{\delta_j},
\end{equation}
where $q_j$ is \eqref{eq:q} for the pair $(L,j)$ over the $H_r$ planning hours, with $\Phi$ evaluated at the probed operating point. The active set takes rivals in order of $E_j$ while the union-bound tail $\sum_{j\notin A}\tau^{-1/2}e^{-E_j/2}$ exceeds $\delta/2$ (at most 3 rivals), skipping rivals whose reach (the largest $q_j$ any dictionary plan could achieve today) is below $\eta_j$. If the skipped rivals' mass exceeds $\delta/2$, or the tail stays at or above $\delta$, the planner buys nothing. The allocation is $\delta_j=(\delta-\text{tail})/|A|$ with $\delta=0.1$. The block is the cheapest dictionary block meeting \eqref{eq:v3}, found by bisection on the level. With banks of up to 150 candidates the tail could not be covered by three rivals, so the planner bought nothing on 862 triggered gates (DEVELOPMENT\_RUNS \S12).

\subsection{Certified planner (\code{cert\_P})}
\begin{lemma}[Certified miss bound; PLANNER\_REVISION Theorem~1]\label{lem:cert}
Assume that under $H_L$ the next $H_r$ hours of measurements are Gaussian around $\mu_L(\bu)$ with covariance $\bar S_{Lj}$, and hold the parameters at their current fits. Let $E_j^+=E_j+\beta_j+\Lambda_j$ be the BIC gap at the end of the window, with $\Lambda_j$ the log-likelihood-ratio increment and $\beta_j=\tfrac12(p_j-p_L)\ln(N^+/N)$ the drift of the penalty. Then
\begin{equation}\label{eq:B}
  \Pr\Big(\sum_{j\in R}e^{-E_j^+}>\tau\ \Big|\ H_L\Big)\le B(\bu)=\frac{1}{\sqrt\tau}\sum_{j\in R}\exp\!\Big(-\frac{E_j+\beta_j}{2}-\frac{q_j(\bu)}{8}\Big).
\end{equation}
\end{lemma}
\begin{proof}
$\Lambda_j\sim\mathcal N(q_j/2,q_j)$, so $\mathbb E[e^{-\Lambda_j/2}]=e^{-q_j/8}$. Square roots are subadditive, so $\sum_j x_j>\tau$ implies $\sum_j\sqrt{x_j}>\sqrt\tau$; Markov's inequality on $\sum_j e^{-E_j^+/2}$ gives \eqref{eq:B}. No independence between rivals is needed.
\end{proof}
The certified planner holds a block for the $H_r$ planning hours and solves
\begin{equation}\label{eq:certP}
  \bu^\star\in\arg\min_{\bu\in\cU}\ c(\bu)+V\,\big[B(\bu)-\delta\big]_+ ,
\end{equation}
over the dictionary for the three nearest rivals and the 24 levels, with the 20 nearest rivals in the sum and the others held at $q_j=0$ (conservative). If some block reaches $B\le\delta$, the hinge is an exact penalty and the planner buys the cheapest such block; otherwise it trades cost against $B$. The v3 condition \eqref{eq:v3} is one fixed allocation of the same budget $\delta$ (PLANNER\_REVISION Corollary~1).

\subsection{Strong greedy (\code{bgreedy\_P}): one-step Bayes-risk minimisation}\label{sec:bgreedy}
\textbf{Classes.} The leader's class is $\cQ_L$. The other candidates with $E_k<25$ are taken in order of $E_k$ and grouped around representatives by the same voltage test as in the definition of equivalence, up to $M=10$ classes. A candidate equivalent to no kept representative opens a new class while fewer than $M$ exist; otherwise it is left out. The class weights are
\begin{equation}\label{eq:W}
  W_c=\sum_{k\in c}w_k\ \ \text{(renormalised over the kept classes)},
\end{equation}
and the representative $r(c)$ (the first, most probable member) supplies the class's predictions.

\begin{lemma}[Bhattacharyya bound on the class-level risk]\label{lem:U}
Treat each class as one hypothesis with prior $W_c$ and Gaussian predictive density $p_c$ for the next hour's data, and let $\hat c$ be the MAP class after that hour. Then
\begin{equation}\label{eq:U}
  \Pr(\hat c\neq c_{\rm true})\le U(\bu)=\sum_{c<c'}\sqrt{W_cW_{c'}}\;\exp\!\Big(-\frac{q_{cc'}(\bu)}{8}\Big),
\end{equation}
with $q_{cc'}$ the separation \eqref{eq:q} of the representatives over the next hour.
\end{lemma}
\begin{proof}
Let $a_c(y)$ be the posterior of class $c$ and $m(y)$ the MAP class. The error mass at $y$ is $1-a_m\le\sum_{c<c'}\min(a_c,a_{c'})$: every class $c\neq m$ appears in the pair $(c,m)$, where the minimum is $a_c$. Integrating against the predictive density,
\[
  \Pr(\text{error})\le\sum_{c<c'}\int\min(W_cp_c,W_{c'}p_{c'})\,dy\le\sum_{c<c'}\sqrt{W_cW_{c'}}\int\sqrt{p_cp_{c'}}\,dy ,
\]
using $\min(x,z)\le\sqrt{xz}$. For Gaussians with covariances $S_c,S_{c'}$ and average $\bar S$, the Bhattacharyya coefficient is $\exp\!\big(-q/8-\tfrac12\ln\frac{\det\bar S}{\sqrt{\det S_c\det S_{c'}}}\big)\le e^{-q/8}$, because the log-determinant term is non-negative.
\end{proof}

\textbf{The rule.} At each gate where a block may be bought,
\begin{equation}\label{eq:bgreedy}
  \boxed{\ \bu^\star\in\arg\min_{\bu\in\{0\}\cup\{\kappa\bs:\ \bs\in\cD,\ 0<\kappa\le\bar\kappa(\bs)\}}\ c(\bu)+V\,U(\bu)\ }
\end{equation}
over the next hour only. Equivalently, buy block $\bu$ only if the risk it removes is worth its price,
\begin{equation}\label{eq:buy}
  V\big(U(0)-U(\bu)\big)>c(\bu),
\end{equation}
and among such blocks take the one with the largest net value $V(U(0)-U(\bu))-c(\bu)$. If none qualifies, buy nothing. $V$ is the price of certainty (Section~\ref{sec:V}); \code{bgreedy\_P} uses $V=1000$~CNY and \code{bgreedy\_lowV\_P} $V=200$~CNY.

\begin{proposition}[Near-optimality of greedy selection]\label{prop:sub}
Fix a menu of blocks and freeze sensitivities and weights, and suppose separations add over executed blocks, $q_{cc'}(\mathcal A)=q_{cc'}(0)+\sum_{b\in\mathcal A}\delta q_{cc'}(b)$ with $\delta q\ge0$. Then the risk reduction
\[
  F(\mathcal A)=U(\varnothing)-U(\mathcal A)=\sum_{c<c'}\sqrt{W_cW_{c'}}\Big(e^{-q_{cc'}(0)/8}-e^{-q_{cc'}(\mathcal A)/8}\Big)
\]
is a non-negative sum of concave functions of modular functions, hence monotone submodular in the set $\mathcal A$ of executed blocks. Greedy selection is therefore within a factor $1-1/e$ of the best schedule of the same length (Nemhauser, Wolsey and Fisher, 1978).
\end{proposition}
Two approximations separate this from the executed policy: separations of different hours add only approximately (the parameter-uncertainty term couples hours, and the cross term $2q^{(1)}\kappa$ can be negative), and the hourly refit changes $W_c$ and the sensitivities, which makes the policy adaptive.

\subsection{Broad check (\code{check\_then\_t\_opt}, dropped)}
For the first two probing hours it buys a zero-sum block with random signs at half of every terminal's headroom, rescaled so that increases equal reductions, and then hands over to the v3 planner. It was dropped after it raised wrong confirmations and cost.

\subsection{Summary}
\begin{center}
\small
\begin{tabularx}{\textwidth}{@{}l>{\raggedright\arraybackslash}X>{\raggedright\arraybackslash}X>{\raggedright\arraybackslash}X>{\raggedright\arraybackslash}X@{}}
\toprule
Strategy & Objective & Rivals considered & Noise model & When it stops buying\\
\midrule
Passive & none & -- & -- & never buys\\
Random & none (random home and sign) & -- & -- & replays the certified planner's hours\\
Crude greedy & $\max\|\Delta\mu_{Lj_1}\|^2$ at the ceiling & nearest rival & white, known parameters & never (only the shared rules)\\
Committee & $\max\sum\bar w_k\|\mu_k-\bar\mu\|^2$ at the ceiling & 8 most probable candidates & white, known parameters & never (only the shared rules)\\
v3 planner & cheapest block with $q_j\ge\eta_j$ & active set of at most 3 & GLS, parameter uncertainty & when no active set covers the tail (almost always)\\
Certified & $\min c+V[B-\delta]_+$ over 2~h & 20 nearest, leader-centred & GLS, parameter uncertainty & when no block lowers $B-\delta$ by more than its cost$/V$\\
Strong greedy & $\min c+V\,U$ over 1~h & all class pairs, up to 10 classes & GLS, parameter uncertainty & when no block lowers $U$ by more than its cost$/V$\\
\bottomrule
\end{tabularx}
\end{center}

\noindent Hard-case test (DEVELOPMENT\_RUNS \S16, 550 ambiguous runs per arm, DR prices): correct confirmations at the end of the day were passive 29\%, random 31\%, committee 37\%, strong greedy 40\% at $V=200$ and at $V=1000$, certified 41\%; daily cost 0, 13.8, 16.3, 16.8, 23.8 and 26.7~CNY. These figures predate the battery-state fix (blocks were sized against the charge before the current block ran) and are optimistic: with the fix the strong greedy is 37--38\% correct at $V=200$--1000 (Section~\ref{sec:V}).

% =====================================================================
\section{The price of certainty \texorpdfstring{$V$}{V}}\label{sec:V}

\textbf{What it means.} $V$ converts misidentification risk into money. By \eqref{eq:buy}, the strong greedy pays at most $V/100$~CNY to lower the bound $U$ by one percentage point: 2~CNY at $V=200$, 10~CNY at $V=1000$. A block costing 4~CNY must therefore lower $U$ by at least 0.02 at $V=200$, or by at least 0.004 at $V=1000$.

\textbf{Two readings.}
\begin{itemize}
  \item \emph{Lagrangian.} $\min c(\bu)+V\,U(\bu)$ is the Lagrangian of ``minimise the risk bound subject to a spend limit''. Every $V$ gives a point on the cost--risk trade-off curve, and choosing $V$ means choosing the point.
  \item \emph{Decision theory.} If naming the wrong class costs the operator $\mathcal L$~CNY (a field visit, or a wrong hosting-capacity decision) and $U$ were the exact error probability, the expected-loss-minimising choice would be $V=\mathcal L$. Because $U$ is an upper bound, $U=k\cdot\Pr(\text{error})$ with $k\ge1$, the consistent setting is $V=\mathcal L/k$.
\end{itemize}

\textbf{How loose is $U$?} The logged hard-case runs of both strong greedy arms record $U(\bu^\star)$ at 3{,}928 planned gates. Table~\ref{tab:calib} compares it with the realised rate of a wrong MAP class at the next gate, when the block's data have arrived.

\begin{table}[h]
\centering
\caption{The risk bound $U$ after the chosen block against the realised wrong-class rate one hour later (DEVELOPMENT\_RUNS \S16 logs, both strong greedy arms).}\label{tab:calib}
\begin{tabular}{@{}lrrrr@{}}
\toprule
$U(\bu^\star)$ & Gates & Mean $U$ & Realised wrong-class rate & Ratio $k$\\
\midrule
0.1--0.2 & 147 & 0.171 & 0.143 & 1.2\\
0.2--0.4 & 1304 & 0.306 & 0.252 & 1.2\\
0.4--0.8 & 1481 & 0.559 & 0.318 & 1.8\\
0.8--1.6 & 839 & 1.096 & 0.434 & 2.5\\
$\ge1.6$ & 138 & 1.994 & 0.645 & 3.1\\
\midrule
All & 3928 & 0.623 & 0.325 & 1.9\\
\bottomrule
\end{tabular}
\end{table}

The bound holds in every bin and is loose by a factor of about $k\approx1.2$--$3$, about 2 overall. With a field visit at $\mathcal L\approx1000$~CNY, the decision-theoretic reading suggests $V\approx500$~CNY.

\textbf{Were 200 and 1000 set by hand?} Yes. $V=1000$~CNY was set to about one field verification visit, and $V=200$~CNY to one fifth of it as a frugal variant. Neither was derived from data. The hard-case test found them equally accurate (9 runs to 9) with $V=200$ cheaper by 30\%.

\textbf{The limit $V\to\infty$.} As $V$ grows, the buy rule \eqref{eq:buy} passes whenever the block lowers $U$ at all, and the minimiser of $c(\bu)+V\,U(\bu)$ moves to the minimiser of $U$ over the admissible set, which sits at the ceiling $\bar\kappa$. A large $V$ therefore cannot buy more hours than the gates allow; it can only make each block larger, up to the ceilings and the caps. A small $V$ buys only the hours where the bound falls steeply.

\textbf{The sweep (DEVELOPMENT\_RUNS \S17).} The strong greedy at six values of $V$, with the committed-block fix, on the 550 ambiguous hard-case runs (Table~\ref{tab:vsweep}). The net value per run is
\begin{equation}
  \mathrm{NV}(V)=\mathcal L\,\frac{C(V)-C_0}{n}-\bar c(V),\qquad \mathcal L=1000~\text{CNY},
\end{equation}
where $C(V)$ and $C_0$ count the runs that confirm the true topology with $V$ and passively, $n=550$, and $\bar c(V)$ is the mean cost a day. The rule, registered before the results were read, takes the largest $\mathrm{NV}$; the secondary rule takes the cheapest $V$ not significantly less accurate than the most accurate one.

\begin{table}[h]
\centering
\caption{The strong greedy against $V$ on the 550 ambiguous hard-case runs (DEVELOPMENT\_RUNS \S17). ``Gained / lost'': runs only this arm, or only passive, got right (exact McNemar, paired within condition).}\label{tab:vsweep}
\begin{tabular}{@{}rrrrrr@{}}
\toprule
$V$ (CNY) & Correct & Gained / lost & Probed h/day & CNY/day & NV (CNY/run)\\
\midrule
passive & 162 & -- & 0 & 0 & 0\\
50 & 190 & 33 / 5 & 1.9 & 6.8 & 44.1\\
100 & 205 & 47 / 4 & 2.6 & 11.6 & 66.6\\
200 & 206 & 47 / 3 & 3.0 & 16.5 & 63.5\\
500 & 210 & 52 / 4 & 3.2 & 21.3 & 66.0\\
1000 & 209 & 52 / 5 & 3.2 & 23.4 & 62.1\\
2000 & 215 & 56 / 3 & 3.2 & 24.7 & 71.6\\
100000 (limit) & 214 & 55 / 3 & 3.2 & 26.4 & 68.2\\
\bottomrule
\end{tabular}
\end{table}

\textbf{Which $V$ is best?} The data do not single one out.
\begin{itemize}
  \item Accuracy rises from $V=50$ to $V=100$ and is then flat within noise up to 2{,}000 ($V=100$ against $V=2000$: 16 runs to 26, $p=0.16$). Probed hours saturate at about 3.2 a day from $V\approx500$, as the limit above predicts; beyond that a larger $V$ only buys larger blocks. The point $V=100{,}000$, run afterwards and never a candidate, confirms the limit: the same accuracy as $V=2000$ (the two differ in 5 runs), the same hours, slightly larger blocks and a lower $\mathrm{NV}$.
  \item $\mathrm{NV}$ is 62--72~CNY per run for every $V$ from 100 to 2{,}000. The registered rule picks $V=2000$ by 5~CNY per run over $V=100$, about 10 correct runs in 550, which is inside the noise. The secondary rule picks $V=100$, at less than half the cost (11.6 against 24.7~CNY a day).
  \item The decision-theoretic value $V\approx\mathcal L/k\approx500$~CNY lies inside the flat range, so the two readings agree.
\end{itemize}
On the sweep's inputs, any $V$ from about 100 to 2{,}000~CNY seemed to give the same value per run. The fresh seeds corrected this (below): $V=2000$ beat $V=100$ by 38 runs to 14 ($p=0.001$), in the same direction as the sweep's 26 to 16. ``Not significantly worse'' was not equivalence.

\textbf{Fresh-seed confirmation (DEVELOPMENT\_RUNS \S18).} On 532 new ambiguous runs (seeds 65--96), the strong greedy at $V=2000$ confirmed the true topology in 43\% of runs against 32\% passive (61 gained, 2 lost), beat random blocks 47 to 8, matched the certified planner (21 to 15) and had the largest net value of the probing arms (88~CNY per run, against 76 certified and 73 committee). On the same inputs $V=100$ was 39\% correct at 10.7~CNY a day (55~CNY per run), against 43\% at 23.2~CNY a day for $V=2000$ (38 runs to 14): the larger $V$ probes more hours (3.2 against 2.5 a day) with larger blocks. With the truth removed from the same banks, it confirmed a wrong topology in 64\% of runs against 60\% passive (37 to 17, $p=0.009$), so the registered safety test failed. The extra wrong confirmations were near-twins of the truth, at most 0.14\% off in voltage on the short lines and 0.27\% under meter stress. The registered decision therefore required a stronger trust gate before the lock. A first check, in which every candidate predicts each hour's meter readings before seeing them, was correctly scaled but did not separate truth-missing banks (areas under the ROC curve 0.47--0.58): the twin's excess misfit over a day is smaller than the random spread of a test of the leader alone. On the 62 inputs whose truth the bank missed by itself (DEVELOPMENT\_RUNS \S19), probing barely changed wrong confirmations (45\% against 44\%, 4 runs to 2): in 51 of them the bank's closure had hit a cap, so the trust gate forbade probing. All of the truth-removed test's harm came from completed banks, a case that arose naturally in 2.4\% of the inputs where probing is allowed.

\textbf{The lock (2026-11-23).} The strong greedy at $V=2000$~CNY, with linear DR settlement and GLS scoring, was locked by the authors' decision, departing from the registered rule. The failed safety test is disclosed with its estimated size: weighted by the 2.4\% rate, the extra wrong confirmations are about 0.1 percentage point of probe-eligible runs (0.4 under meter stress), against 11 points of extra correct confirmations, and each extra wrong topology was within 0.27\% of the truth in voltage. Stage~M keeps measuring truth-removed and natural truth-missing runs.

% =====================================================================
\appendix
\section{Code}\label{app:code}
The listings below are copied from the study's source files (\code{rnj\_wzzt\_dr/experiments/hourly\_probing\_20261118/}) when this file is built; the line numbers are those of the files.

\subsection{GLS whitening, offsets and the hourly refit}

\begin{lstlisting}[firstnumber=50,title={\code{model.py}, \code{root\_gamma}, lines 50--61}]
def root_gamma(V0sq, Y):
    """Whitening strength for the shared root-meter error.

    Every terminal's y = V0^2 - V_i^2 subtracts the same root reading, so the
    noise of one time slot has covariance c (I + rho 11'), rho = V0^4 / V^4.
    Its inverse square root is I - gamma 11'/n with gamma = 1 - 1/sqrt(1 + n rho).
    """
    Y = np.asarray(Y, float)
    n = Y.shape[1]
    v0sq = np.asarray(V0sq, float).reshape(-1)
    rho = float(np.mean(v0sq**2)/np.mean((v0sq[:, None]-Y)**2))
    return 1.-1./np.sqrt(1.+n*rho)
\end{lstlisting}

\begin{lstlisting}[firstnumber=64,title={\code{model.py}, \code{whiten\_slots}, lines 64--71}]
def whiten_slots(A, n, gamma):
    """Apply I - gamma 11'/n to every time slot of rows ordered (t, i) (GLS whitening)."""
    A = np.asarray(A, float)
    if not gamma:
        return A
    T = A.shape[0]//n
    B = A.reshape(T, n, -1)
    return (B-gamma*B.mean(axis=1, keepdims=True)).reshape(A.shape)
\end{lstlisting}

\begin{lstlisting}[firstnumber=37,title={\code{model.py}, \code{demean\_rows}, lines 37--47}]
def demean_rows(A, n):
    """Per-terminal demeaning of rows ordered (t, i).

    Fitting on demeaned rows is the same as fitting a free offset per meter
    (Frisch-Waugh-Lovell), also when the other coefficients are bounded: for any
    fixed theta the best offset is the terminal's mean residual.
    """
    A = np.asarray(A, float)
    T = A.shape[0]//n
    B = A.reshape(T, n, -1)
    return (B-B.mean(axis=0, keepdims=True)).reshape(A.shape)
\end{lstlisting}

\begin{lstlisting}[firstnumber=162,title={\code{model.py}, \code{Fit.refit}, lines 162--196}]
    def refit(self, Phi, y, lam, steps=1, P=None, Q=None, v0sq=None, use_loss=True, offsets=False, gamma=0.):
        """Bounded LS on y - lambda(theta-); `steps` fixed-point steps of the loss term.

        With `offsets`, every meter gets a free constant (fitted by demeaning);
        theta stays the physical model and `self.offset` holds the meter offsets.
        With `gamma` > 0 the rows are GLS-whitened for the shared root-meter error
        (`whiten_slots`); whitening and demeaning commute.
        """
        n = self.c.n
        A = whiten_slots(Phi, n, gamma)
        A = demean_rows(A, n) if offsets else A
        for step in range(steps):
            target = y-(lam if lam is not None else 0.)
            if not np.all(np.isfinite(target)):
                self.failures = getattr(self, "failures", 0)+1
                target = y
            b = whiten_slots(target[:, None], n, gamma).ravel()
            b = demean_rows(b[:, None], n).ravel() if offsets else b
            try:
                sol = lsq_linear(A, b, bounds=(self.lb, self.ub), method="bvls", tol=1e-12,
                                 max_iter=5000)
            except np.linalg.LinAlgError:
                self.failures = getattr(self, "failures", 0)+1
                sol = lsq_linear(A, b, bounds=(self.lb, self.ub), method="trf", tol=1e-12,
                                 max_iter=5000)
            self.theta = np.clip(sol.x, self.lb, self.ub)
            if use_loss and step < steps-1:
                lam = self.c.loss_term(self.theta, P, Q, v0sq).ravel()
        resid = b-A@self.theta
        self.sse = float(resid@resid)
        self.offset = ((target-Phi@self.theta).reshape(-1, n).mean(axis=0) if offsets else np.zeros(n))
        tol = 1e-9
        self.free = (self.theta > self.lb+tol) & (self.theta < self.ub-tol)
        self.G = A.T@A
        return self
\end{lstlisting}

\begin{lstlisting}[firstnumber=206,title={\code{model.py}, \code{Fit.covariance}, lines 206--219}]
    def covariance(self, sigma2):
        """Sigma_k on free coordinates (18); bound coordinates get zero."""
        p = len(self.theta)
        cov = np.zeros((p, p))
        f = np.flatnonzero(self.free)
        if len(f):
            G = self.G[np.ix_(f, f)]
            G = G+1e-12*np.trace(G)/len(f)*np.eye(len(f))
            try:
                inv = cho_solve(cho_factor(G), np.eye(len(f)))
            except np.linalg.LinAlgError:
                inv = np.linalg.pinv(G)
            cov[np.ix_(f, f)] = sigma2*inv
        return cov
\end{lstlisting}

\subsection{Response model: design and operating-point sensitivity}

\begin{lstlisting}[firstnumber=99,title={\code{model.py}, \code{Candidate.design}, lines 99--105}]
    def design(self, P, Q):
        """Phi with rows (t, i) t-major and columns (r_C..., x_C...)."""
        a, b = self.flows(P, Q)
        T = P.shape[0]
        Fr = (a[:, None, :]*self.Z.T[None, :, :]).reshape(T*self.n, self.S)
        Fx = (b[:, None, :]*self.Z.T[None, :, :]).reshape(T*self.n, self.S)
        return np.hstack((Fr, Fx))
\end{lstlisting}

\begin{lstlisting}[firstnumber=131,title={\code{model.py}, \code{Candidate.sensitivity}, lines 131--142}]
    def sensitivity(self, theta, P, Q, v0sq, loss=True, h=1e-5):
        """Operating-point sensitivity dy/dp per sample: T x n x n (R_k + J_k)."""
        r = theta[:self.S]
        R = (self.Z.T*r)@self.Z
        T = P.shape[0]
        M = np.repeat(R[None], T, axis=0)
        if loss:
            for m in range(self.n):
                Pp = P.copy(); Pp[:, m] += h
                Pm = P.copy(); Pm[:, m] -= h
                M[:, :, m] += (self.loss_term(theta, Pp, Q, v0sq)-self.loss_term(theta, Pm, Q, v0sq))/(2*h)
        return M
\end{lstlisting}

\subsection{Scores, equivalence and confirmation}

\begin{lstlisting}[firstnumber=8,title={\code{scores.py}, \code{bic}, lines 8--12}]
def bic(sses, ps, N):
    sses, ps = np.asarray(sses, float), np.asarray(ps, float)
    sigma2 = float(np.min(sses/(N-ps)))
    J = sses/(2*sigma2)+.5*ps*math.log(N)
    return J, sigma2
\end{lstlisting}

\begin{lstlisting}[firstnumber=31,title={\code{scores.py}, \code{posterior\_ok}, lines 31--34}]
def posterior_ok(E, equivalent, leader, alpha):
    rivals = [k for k in range(len(E)) if k != leader and not equivalent[k]]
    total = float(np.sum(np.exp(-np.asarray(E)[rivals]))) if rivals else 0.
    return total <= alpha/(1-alpha), total
\end{lstlisting}

\begin{lstlisting}[firstnumber=145,title={\code{run.py}, \code{excerpt}, lines 145--174}]
        N = y.size
        J, sigma2 = scores.bic([f.sse for f in fits], [c.p+(n if s.meter_offsets else 0) for c in cands], N)
        L = int(np.argmin(J)); E = J-J[L]
        sigma2_prev = sigma2
        w = np.exp(-(J-J.min())); w /= w.sum()
        thetas = [f.theta for f in fits]
        near = [k for k in range(K) if k != L and E[k] < 25]
        eq = np.zeros(K, bool); eq[L] = True
        if near:
            vL = scores.voltages(cands[L], thetas[L], Pf, Qf, V0f, s.use_loss_term)
            for k in near:
                vk = scores.voltages(cands[k], thetas[k], Pf, Qf, V0f, s.use_loss_term)
                eq[k] = float(np.max(np.abs(vk-vL))) <= s.equivalence_pu
        ok, mass = scores.posterior_ok(E, eq, L, s.alpha)
        if arm == "check_then_t_opt" and g <= status["check_until"]:
            ok = False
        if ok and status["leader"] == L:
            status["count"] += 1
        elif ok:
            status["leader"], status["count"] = L, 1
        else:
            status["leader"], status["count"] = None, 0
        was = status["confirmed"]
        status["confirmed"] = status["count"] >= s.confirm_hours
        if status["confirmed"] and cands[L].key != inp.truth_key:
            ever_wrong = True; ever_wrong_types.add(scores.rival_type(cands[L].key, inp.truth_key))
            wrong_leaders.add(L)
        if status["confirmed"] and cands[L].key == inp.truth_key and first_correct is None:
            first_correct = g
        rivals = sorted([k for k in range(K) if k != L and not eq[k]], key=lambda k: E[k])
\end{lstlisting}

\subsection{When a block may be bought; dispatch to the strategies; no-progress rule}

\begin{lstlisting}[firstnumber=225,title={\code{run.py}, \code{excerpt}, lines 225--257}]
            elif probing and not ok and not status["confirmed"] and not status["stopped"] and rivals \
                    and (not s.trust_gate or arm in ("t_opt_P", "check_then_t_opt")
                         or (closure_ok and gate["adequacy"] <= s.adequacy_max)):
                if arm == "check_then_t_opt" and status["check_done"] < s.check_hours:
                    u = planner.check_block(s, inp.res, hb, soc_c, rng_chk)
                    status["check_done"] += 1; status["check_until"] = hb+1
                    plan = dict(reason="check", shape="zero_sum_check", active=[], beyond=[], reach={},
                                cost=planner.block_cost(s, inp.res, hb, u), window_empty=False)
                else:
                    covs = [f.covariance(sigma2) for f in fits]
                    gt = planner.Gate(inp, s, cands, thetas, covs, sigma2, E, L, eq, w, (Pf, Qf, V0f),
                                      soc_c, spent_c, hb)
                    if s.meter_offsets:
                        gt.offsets = [f.offset for f in fits]
                    gt.N_rows = N
                    gt.root_gamma = gamma
                    if arm == "cert_P":
                        plan = gt.plan_certified(rivals)
                    elif arm.startswith("bgreedy"):
                        plan = gt.plan_bayes_greedy(rivals)
                    elif arm in ("qbc_P", "greedy_P"):
                        plan = gt.plan_heuristic(rivals, arm[:-2])
                    else:
                        plan = gt.plan(rivals)
                    u = plan["block_kw"]
                    last_beyond = list(plan["beyond"]); last_reach = plan["reach"]
                    if plan["active"]:
                        G = min(E[j] for j in plan["active"])
                        if np.any(u != 0):
                            progress.append(G)
                            if len(progress) > s.no_progress_hours and \
                                    progress[-1]-progress[-1-s.no_progress_hours] < s.no_progress_gain:
                                status["stopped"] = True
\end{lstlisting}

\subsection{Committed spend and battery state}

\begin{lstlisting}[firstnumber=55,title={\code{run.py}, \code{committed}, lines 55--66}]
def committed(s, res, g, hb, blocks, spent, soc):
    """Spend and battery state after the blocks already sealed for hours g..hb-1 but not yet run.

    The block for hour hb is sealed at gate g while the block for hour g is still to run, so the
    daily budget and the battery energy must count it (fixed 2026-11-23, DEVELOPMENT_RUNS section 16).
    """
    soc = np.array(soc, float)
    for h in range(g, min(hb, s.hours)):
        if np.any(blocks[h] != 0):
            spent += planner.block_cost(s, res, h, blocks[h])
            soc = sc.update_soc(s, res, soc, sc.split_block(s, res, h, soc, blocks[h]))
    return float(spent), soc
\end{lstlisting}

\subsection{Prediction rows and GLS in the planner}

\begin{lstlisting}[firstnumber=68,title={\code{planner.py}, \code{Gate.rows}, lines 68--71}]
    def rows(self, k, hours):
        d = self.cdata(k)
        idx = np.concatenate([np.arange(SPH*(h-self.hb), SPH*(h-self.hb)+SPH) for h in hours])
        return d["mu"][idx], d["M"][idx], d["slots"][idx]
\end{lstlisting}

\begin{lstlisting}[firstnumber=60,title={\code{planner.py}, \code{Gate.wh}, lines 60--63}]
    def wh(self, X):
        """GLS-whiten a T x n array per time slot (shared root-meter error); identity if gamma = 0."""
        g = getattr(self, "root_gamma", 0.)
        return X-g*X.mean(axis=1, keepdims=True) if g else X
\end{lstlisting}

\subsection{Dictionary of shapes}

\begin{lstlisting}[firstnumber=113,title={\code{planner.py}, \code{Gate.shapes}, lines 113--125}]
    def shapes(self, rivals):
        n, out = self.n, []
        for i in range(n):
            for sgn in (1., -1.):
                s = np.zeros(n); s[i] = sgn
                out.append(("node%d%+d" % (i, int(sgn)), s))
        for j in rivals:
            nodes = self.differing_nodes(j)
            if 1 < len(nodes) < n:
                for sgn in (1., -1.):
                    s = np.zeros(n); s[nodes] = sgn
                    out.append(("group%d%+d" % (j, int(sgn)), s))
        return out
\end{lstlisting}

\begin{lstlisting}[firstnumber=127,title={\code{planner.py}, \code{Gate.differing\_nodes}, lines 127--135}]
    def differing_nodes(self, j):
        fl, fj = self.cands[self.L].family, self.cands[j].family
        nodes = []
        for i in range(self.n):
            cl = frozenset(C for C in fl if i in C and len(C) > 1)
            cj = frozenset(C for C in fj if i in C and len(C) > 1)
            if cl != cj:
                nodes.append(i)
        return nodes
\end{lstlisting}

\subsection{Ceilings, budget and cost}

\begin{lstlisting}[firstnumber=138,title={\code{planner.py}, \code{Gate.headroom\_hour}, lines 138--148}]
    def headroom_hour(self, h, hours_held=1, energy=True):
        res, s = self.res, self.s
        lo, hi = s.soc_bounds
        e = res["battery_kwh"]
        if energy:
            up_b = np.minimum(res["battery_kw"], np.maximum((hi-self.soc)*e/(s.battery_eta*hours_held), 0.))
            dn_b = np.minimum(res["battery_kw"], np.maximum((self.soc-lo)*e*s.battery_eta/hours_held, 0.))
        else:
            up_b = dn_b = res["battery_kw"]
        ev = res["ev_kw"] if sc.ev_available(s, h) else 0.*res["ev_kw"]
        return up_b+ev, dn_b+s.ev_down_fraction*ev
\end{lstlisting}

\begin{lstlisting}[firstnumber=150,title={\code{planner.py}, \code{Gate.cexp}, lines 150--153}]
    @property
    def cexp(self):
        """Exponent a of the cost in the level: c(k s) = k^a c(s); 2 for contract (3), 1 for linear DR."""
        return 1 if self.s.cost_model == "linear" else 2
\end{lstlisting}

\begin{lstlisting}[firstnumber=155,title={\code{planner.py}, \code{Gate.level\_cost}, lines 155--156}]
    def level_cost(self, k, unit):
        return float(unit*k**self.cexp)
\end{lstlisting}

\begin{lstlisting}[firstnumber=158,title={\code{planner.py}, \code{Gate.level\_for\_cost}, lines 158--159}]
    def level_for_cost(self, cost, unit):
        return (max(cost, 0.)/unit)**(1./self.cexp)
\end{lstlisting}

\begin{lstlisting}[firstnumber=161,title={\code{planner.py}, \code{Gate.unit\_cost}, lines 161--170}]
    def unit_cost(self, s_kw, h):
        """Hourly cost (CNY) of the block s_kw at unit level: contract (3) or linear DR settlement."""
        up, dn = self.headroom_hour(h, energy=False)
        cap = np.where(s_kw >= 0, up, dn)
        if np.any((cap <= 1e-9) & (np.abs(s_kw) > 0)):
            return np.inf
        if self.s.cost_model == "linear":
            return float(np.sum(np.where(s_kw >= 0, self.s.price_up_cny, self.s.price_down_cny)*np.abs(s_kw)))
        price = sc.node_price(self.s, self.res, h)
        return float(np.sum(price*s_kw**2/np.maximum(cap, 1e-9)))
\end{lstlisting}

\begin{lstlisting}[firstnumber=172,title={\code{planner.py}, \code{Gate.ceilings\_hour}, lines 172--236}]
    def ceilings_hour(self, s_kw, h, energy=True, hours_held=1):
        """kappa caps for one hour: headroom/energy, voltage, transformer, unbalance, linearisation."""
        s, n = self.s, self.n
        up, dn = self.headroom_hour(h, hours_held, energy)
        cap = np.where(s_kw >= 0, up, dn)
        nz = np.abs(s_kw) > 1e-12
        kA = float(np.min(cap[nz]/np.abs(s_kw[nz]))) if np.any(nz) else np.inf
        s_pu = s_kw/self.base_kw
        sl = hour_slots(h)
        V0 = self.V0f[sl]
        # robust voltage over plausible candidates
        kV = np.inf
        lo_band, hi_band = s.voltage_band
        for k in self.plausible:
            mu, M, _ = self.rows(k, [h])
            c, th = self.cands[k], self.thetas[k]
            vkey = ("volt", k, h)
            if vkey not in self._cache:
                loss = s.use_loss_term
                f_lo, f_hi = 1-2*s.forecast_sigma, 1+2*s.forecast_sigma
                self._cache[vkey] = (
                    np.sqrt(np.maximum(V0[:, None]-mu, .25)),
                    np.sqrt(np.maximum(V0[:, None]-c.predict(th, self.Pf[sl]*f_lo, self.Qf[sl]*f_lo, V0, loss), .25)),
                    np.sqrt(np.maximum(V0[:, None]-c.predict(th, self.Pf[sl]*f_hi, self.Qf[sl]*f_hi, V0, loss), .25)))
            v_hat, hv, lv = self._cache[vkey]
            g1 = -np.einsum("tim,m->ti", M, s_pu)/(2*v_hat)
            flow = c.Z@s_pu                                  # probe flow per support
            phi = c.Z.T*flow                                 # n x S (r columns)
            cov_r = self.covs[k][:c.S, :c.S]
            b1 = np.sqrt(np.maximum(np.einsum("js,st,jt->j", phi, cov_r, phi), 0.))/(2*v_hat)
            up_den = g1+s.z_voltage*b1
            dn_den = -g1+s.z_voltage*b1
            with np.errstate(divide="ignore", invalid="ignore"):
                a = np.where(up_den > 0, np.maximum(hi_band-hv, 0.)/up_den, np.inf)
                b = np.where(dn_den > 0, np.maximum(lv-lo_band, 0.)/dn_den, np.inf)
            kV = min(kV, float(np.min(a)), float(np.min(b)))
        # transformer
        P0 = self.Pf[sl].sum(axis=1); Q0 = self.Qf[sl].sum(axis=1)
        dP = float(s_pu.sum())
        room = np.sqrt(np.maximum(self.smax**2-Q0**2, 0.))
        if abs(dP) < 1e-15:
            kS = np.inf
        elif dP > 0:
            kS = float(np.min(np.maximum(room-P0, 0.)/dP))
        else:
            kS = float(np.min(np.maximum(room+P0, 0.)/(-dP)))
        # linearisation remainder (27), leader parameters
        c, th = self.cands[self.L], self.thetas[self.L]
        r, x = th[:c.S], th[c.S:]
        a_f, b_f = c.flows(self.Pf[sl], self.Qf[sl])
        flow_bg = np.hypot(a_f, b_f)                        # T x S
        xi = 2*s.forecast_sigma*flow_bg
        dS = np.abs(c.Z@s_pu)                               # S
        muL, _, _ = self.rows(self.L, [h])
        v2 = np.maximum(V0[:, None]-muL, .25)               # T x n
        w = .25*(r**2+x**2)                                 # S
        Aq = (w*dS**2)@c.Z                                  # n
        Bq = (2*w*dS*xi)@c.Z                                # T x n
        sig_e = math.sqrt(self.sigma2)
        Aq = Aq[None, :]/v2; Bq = Bq/v2
        with np.errstate(divide="ignore", invalid="ignore"):
            root = np.where(Aq > 0, (-Bq+np.sqrt(Bq**2+4*Aq*sig_e))/(2*Aq),
                            np.where(Bq > 0, sig_e/Bq, np.inf))
        kL = float(np.min(root))
        return dict(A=kA, V=kV, S=kS, U=kV/s.unbalance_factor, L=kL)
\end{lstlisting}

\begin{lstlisting}[firstnumber=238,title={\code{planner.py}, \code{Gate.window}, lines 238--253}]
    def window(self, s_kw):
        parts = dict(A=np.inf, V=np.inf, S=np.inf, U=np.inf, L=np.inf)
        for h in self.plan_hours:
            c = self.ceilings_hour(s_kw, h, energy=True, hours_held=self.tr)
            for key in parts:
                parts[key] = min(parts[key], c[key])
        unit = max(self.unit_cost(s_kw, h) for h in self.plan_hours)
        if not np.isfinite(unit):
            parts["B"] = 0.
        elif unit <= 0:
            parts["B"] = np.inf
        else:
            left = max(self.s.daily_cap_cny-self.spent, 0.)
            parts["B"] = min(self.level_for_cost(self.s.hourly_cap_cny, unit),
                             self.level_for_cost(left/self.tr, unit))
        return min(parts.values()), parts, unit
\end{lstlisting}

\begin{lstlisting}[firstnumber=581,title={\code{planner.py}, \code{block\_cost}, lines 581--591}]
def block_cost(settings, res, hour, u_kw):
    """Cost of an executed block, one hour: contract (3) or linear DR settlement per kWh moved."""
    u_kw = np.asarray(u_kw, float)
    if settings.cost_model == "linear":
        return float(settings.price_up_cny*np.sum(np.maximum(u_kw, 0.))
                     + settings.price_down_cny*np.sum(np.maximum(-u_kw, 0.)))
    up = res["battery_kw"]+(res["ev_kw"] if sc.ev_available(settings, hour) else 0.)
    dn = res["battery_kw"]+settings.ev_down_fraction*(res["ev_kw"] if sc.ev_available(settings, hour) else 0.)
    cap = np.where(u_kw >= 0, up, dn)
    price = sc.node_price(settings, res, hour)
    return float(np.sum(price*u_kw**2/np.maximum(cap, 1e-9)))
\end{lstlisting}

\subsection{Random block}

\begin{lstlisting}[firstnumber=612,title={\code{planner.py}, \code{random\_block}, lines 612--630}]
def random_block(settings, res, hour, cost, rng, soc):
    """Random single-node block with the given hourly cost, within headroom."""
    n = len(res["battery_kw"])
    i = int(rng.integers(n)); sgn = rng.choice([-1., 1.])
    up = res["battery_kw"][i]+(res["ev_kw"][i] if sc.ev_available(settings, hour) else 0.)   # contract basis only
    # the cap is the state-of-charge-aware headroom every other arm uses (2026-11-23: the old cap let a full
    # battery charge at its rated power, which run.py's headroom_excess record caught in section 18)
    hu, hd = sc.headroom(settings, res, hour, soc)
    cap = hu[i] if sgn > 0 else hd[i]
    if settings.cost_model == "linear":
        unit = settings.price_up_cny if sgn > 0 else settings.price_down_cny
        u = np.zeros(n); u[i] = sgn*min(max(cost, 0.)/unit, cap)
        return u
    price = sc.node_price(settings, res, hour)[i]
    contract = up if sgn > 0 else res["battery_kw"][i]+settings.ev_down_fraction*(
        res["ev_kw"][i] if sc.ev_available(settings, hour) else 0.)
    level = math.sqrt(max(cost, 0.)*contract/price) if price > 0 else 0.
    u = np.zeros(n); u[i] = sgn*min(level, cap)
    return u
\end{lstlisting}

\subsection{Crude greedy and query by committee}

\begin{lstlisting}[firstnumber=472,title={\code{planner.py}, \code{Gate.plan\_heuristic}, lines 472--510}]
    def plan_heuristic(self, rivals_sorted, kind, target_cost=None):
        """Simple rules for comparison: single-node +/- blocks scored on the next hour only.

        greedy: largest |mu_L - mu_j1|^2 against the nearest non-equivalent rival j1.
        qbc:    largest posterior-weighted disagreement sum_k w_k |mu_k - mu_bar|^2 over the
                plausible candidates (query by committee).
        The level is the hour's ceiling (window incl. caps), or the level whose cost equals
        target_cost (matched-cost variant).
        """
        out = dict(block_kw=np.zeros(self.n), reason=kind, active=list(rivals_sorted[:1]), beyond=[], reach={},
                   kappa=0., cost=0., window_empty=False, shape=None, q_planned={}, eta={}, parts=None)
        if self.tr <= 0 or (target_cost is not None and target_cost <= 0):
            return out
        h = self.plan_hours[0]
        if kind == "greedy":
            members, w = [self.L, rivals_sorted[0]], None
        else:
            members = list(self.plausible)
            w = np.array([self.w[k] for k in members]); w = w/w.sum()
        rows = {k: self.rows(k, [h]) for k in members}
        best = (-np.inf, None, None, 0.)
        for name, shape in self.shapes([]):                  # single-node shapes only
            s_kw = shape*self.sig_kw
            kmax, parts, unit = self.window(s_kw)
            if not np.isfinite(kmax) or kmax <= 0 or not np.isfinite(unit) or unit <= 0:
                continue
            k = kmax if target_cost is None else min(kmax, self.level_for_cost(target_cost, unit))
            u = np.repeat((k*s_kw/self.base_kw)[None, :], SPH, axis=0)
            pred = {m: rows[m][0]+np.einsum("tim,tm->ti", rows[m][1], u) for m in members}
            if kind == "greedy":
                score = float(np.sum((pred[members[0]]-pred[members[1]])**2))
            else:
                mbar = sum(wi*pred[m] for wi, m in zip(w, members))
                score = float(sum(wi*np.sum((pred[m]-mbar)**2) for wi, m in zip(w, members)))
            if score > best[0]:
                best = (score, name, k*s_kw, self.level_cost(k, unit))
        if best[1] is not None:
            out.update(shape=best[1], block_kw=best[2], cost=best[3], kappa=float(np.max(np.abs(best[2]))))
        return out
\end{lstlisting}

\subsection{v3 dictionary planner}

\begin{lstlisting}[firstnumber=74,title={\code{planner.py}, \code{Gate.threshold}, lines 74--76}]
    def threshold(self):
        """1/t with t = alpha/(1-alpha); halved threshold (x2) when the tail uses order 1."""
        return (1-self.s.alpha)/self.s.alpha*(2. if self.s.tail_order == 1. else 1.)
\end{lstlisting}

\begin{lstlisting}[firstnumber=78,title={\code{planner.py}, \code{Gate.c\_tail}, lines 78--81}]
    def c_tail(self, j):
        if self.s.tail_order == 1.:
            return self.threshold()*math.exp(-self.E[j])
        return math.sqrt(self.threshold())*math.exp(-.5*self.E[j])
\end{lstlisting}

\begin{lstlisting}[firstnumber=83,title={\code{planner.py}, \code{Gate.gamma}, lines 83--84}]
    def gamma(self, j, delta_j):
        return self.s.omega*(math.log(self.threshold())-self.E[j])+math.log(1./delta_j)
\end{lstlisting}

\begin{lstlisting}[firstnumber=86,title={\code{planner.py}, \code{Gate.eta}, lines 86--87}]
    def eta(self, j, delta_j):
        return max(2*self.gamma(j, delta_j)/(self.s.omega*(1-self.s.omega)), 0.)
\end{lstlisting}

\begin{lstlisting}[firstnumber=89,title={\code{planner.py}, \code{Gate.delta\_mu}, lines 89--93}]
    def delta_mu(self, j, hours, u_pu_hours):
        muL, ML, _ = self.rows(self.L, hours)
        muj, Mj, _ = self.rows(j, hours)
        U = np.repeat(np.asarray(u_pu_hours), SPH, axis=0)          # T x n
        return (muL-muj)+np.einsum("tim,tm->ti", ML-Mj, U)
\end{lstlisting}

\begin{lstlisting}[firstnumber=95,title={\code{planner.py}, \code{Gate.q\_full}, lines 95--106}]
    def q_full(self, j, hours, u_pu_hours):
        """q_j(U) = dmu' S_jw^-1 dmu with parameter uncertainty (15a)."""
        dmu = self.wh(self.delta_mu(j, hours, u_pu_hours)).ravel()
        U = np.repeat(np.asarray(u_pu_hours), SPH, axis=0)
        _, _, sl = self.rows(j, hours)
        P, Q = self.Pf[sl]+U, self.Qf[sl]
        S = self.sigma2*np.eye(dmu.size)
        for k, wgt in ((self.L, self.s.omega), (j, 1-self.s.omega)):
            Phi = self.whp(self.cands[k].design(P, Q))
            S += wgt*Phi@self.covs[k]@Phi.T
        cf = cho_factor(S)
        return float(dmu@cho_solve(cf, dmu))
\end{lstlisting}

\begin{lstlisting}[firstnumber=108,title={\code{planner.py}, \code{Gate.q\_frozen\_hour}, lines 108--110}]
    def q_frozen_hour(self, j, h, u_pu):
        dmu = self.wh(self.delta_mu(j, [h], [u_pu]))
        return float(np.sum(dmu**2)/self.sigma2)
\end{lstlisting}

\begin{lstlisting}[firstnumber=256,title={\code{planner.py}, \code{Gate.reach}, lines 256--271}]
    def reach(self, j):
        """Q_bar_j: frozen-form bound on q_j for any dictionary plan of the rest of the day."""
        total = 0.
        shapes = self.shapes([j])
        for h in self.rest_hours:
            best = self.q_frozen_hour(j, h, np.zeros(self.n))
            for _, s in shapes:
                s_kw = s*self.sig_kw
                key = ("ceil", h, s.tobytes())
                if key not in self._cache:
                    self._cache[key] = min(self.ceilings_hour(s_kw, h, energy=False).values())
                kbar = self._cache[key]
                if np.isfinite(kbar) and kbar > 0:
                    best = max(best, self.q_frozen_hour(j, h, kbar*s_kw/self.base_kw))
            total += best
        return total
\end{lstlisting}

\begin{lstlisting}[firstnumber=274,title={\code{planner.py}, \code{Gate.active\_set}, lines 274--309}]
    def active_set(self, rivals_sorted):
        s = self.s
        c = {j: self.c_tail(j) for j in rivals_sorted}
        total = sum(c.values())
        active, beyond, reach_log = [], [], {}
        tail = total
        if not s.tail_guarantee:
            # best-effort variant (development only): the first max_active reachable
            # rivals by BIC gap, scanning at most 10; the union-bound tail is not required
            for j in rivals_sorted[:10]:
                if len(active) >= s.max_active:
                    break
                Q = self.reach(j)
                eta_lenient = self.eta(j, s.delta)
                reach_log[j] = (Q, eta_lenient)
                (beyond if Q < eta_lenient else active).append(j)
            if not active:
                return [], beyond, reach_log, tail, None
            return active, beyond, reach_log, tail-sum(c[j] for j in active), s.delta/len(active)
        for j in rivals_sorted:
            if tail <= s.delta/2 or len(active) >= s.max_active:
                break
            Q = self.reach(j)
            eta_lenient = self.eta(j, s.delta)
            reach_log[j] = (Q, eta_lenient)
            if Q < eta_lenient:
                beyond.append(j)
            else:
                active.append(j); tail -= c[j]
        beyond_mass = sum(c[j] for j in beyond)
        if beyond_mass > s.delta/2 or tail >= s.delta:
            return [], beyond, reach_log, tail, None
        if not active:
            return [], beyond, reach_log, tail, None
        delta_j = (s.delta-tail)/len(active)
        return active, beyond, reach_log, tail, delta_j
\end{lstlisting}

\begin{lstlisting}[firstnumber=513,title={\code{planner.py}, \code{Gate.plan}, lines 513--578}]
    def plan(self, rivals_sorted):
        s = self.s
        out = dict(block_kw=np.zeros(self.n), reason="", active=[], beyond=[], reach={}, kappa=0.,
                   cost=0., window_empty=False, shape=None, q_planned={}, eta={}, parts=None)
        if self.tr <= 0:
            out["reason"] = "no_hours_left"; return out
        active, beyond, reach_log, tail, delta_j = self.active_set(rivals_sorted)
        out.update(active=active, beyond=beyond, reach={int(k): [float(a), float(b)] for k, (a, b) in reach_log.items()},
                   tail=float(tail))
        if not active:
            out["reason"] = "beyond_reach" if beyond else "no_active"; return out
        etas = {j: self.eta(j, delta_j) for j in active}
        out["eta"] = {int(j): float(v) for j, v in etas.items()}
        active = [j for j in active if etas[j] > 0]          # eta_j = 0: already met
        if not active:
            out["reason"] = "already_met"; return out
        zero = [np.zeros(self.n)]*self.tr
        if all(self.q_full(j, self.plan_hours, zero) >= etas[j] for j in active):
            out["reason"] = "background_suffices"; return out
        best, fallback = None, None
        for name, shape in self.shapes(active):
            s_kw = shape*self.sig_kw
            kmax, parts, unit = self.window(s_kw)
            if not np.isfinite(kmax) or kmax <= 0 or not np.isfinite(unit):
                continue

            def ratio(k):
                u = [k*s_kw/self.base_kw]*self.tr
                return min(self.q_full(j, self.plan_hours, u)/etas[j] for j in active)
            grid = np.linspace(0., kmax, s.kappa_grid+1)[1:]
            prev, hit = 0., None
            for k in grid:
                if ratio(k) >= 1.:
                    hit = k; break
                prev = k
            if hit is None:
                r_top = ratio(kmax)
                if fallback is None or r_top > fallback[0]:
                    fallback = (r_top, name, s_kw, kmax, unit, parts)
                continue
            lo, hi = prev, hit
            for _ in range(s.bisection_steps):
                mid = .5*(lo+hi)
                if ratio(mid) >= 1.:
                    hi = mid
                else:
                    lo = mid
                if hi-lo <= 1e-4*hi:
                    break
            cost = self.level_cost(hi, unit)
            if best is None or cost < best[0]:
                best = (cost, name, s_kw, hi, unit, parts, lo)
        if best is not None:
            cost, name, s_kw, k, unit, parts, below = best
            out.update(reason="admissible", shape=name, kappa=float(k), kappa_below=float(below))
        elif fallback is not None:
            _, name, s_kw, k, unit, parts = fallback
            out.update(reason="window_empty", shape=name, kappa=float(k), window_empty=True)
        else:
            out["reason"] = "no_feasible_shape"; return out
        out["block_kw"] = k*s_kw
        out["cost"] = self.level_cost(k, self.unit_cost(s_kw, self.hb))
        out["parts"] = {key: float(v) for key, v in parts.items()}
        u = [out["block_kw"]/self.base_kw]*self.tr
        out["q_planned"] = {int(j): float(self.q_full(j, self.plan_hours, u)) for j in active}
        return out
\end{lstlisting}

\subsection{Certified planner}

\begin{lstlisting}[firstnumber=312,title={\code{planner.py}, \code{Gate.cert\_coef}, lines 312--319}]
    def cert_coef(self, j, t):
        """t^-1/2 exp(-(E_j + beta_j)/2): rival j's weight in the miss bound B (Theorem 1)."""
        N_now = getattr(self, "N_rows", None)
        beta = 0.
        if N_now:
            N_end = N_now+SPH*self.tr*self.n
            beta = .5*(self.cands[j].p-self.cands[self.L].p)*math.log(N_end/N_now)
        return math.exp(-.5*(self.E[j]+beta))/math.sqrt(t)
\end{lstlisting}

\begin{lstlisting}[firstnumber=321,title={\code{planner.py}, \code{Gate.plan\_certified}, lines 321--377}]
    def plan_certified(self, rivals_sorted):
        """min_u c_h(u) + V [B(u) - delta]_+ over dictionary shapes x levels (PLANNER_REVISION.md section 5).

        B(u) = t^-1/2 sum_j exp(-(E_j+beta_j)/2 - q_j(u)/8) over the cert_rivals nearest rivals,
        the others held at q_j = 0 (conservative). q_j is quadratic in the level of a shape, with
        S_j evaluated at the forecast operating point (one Cholesky factor per rival).
        """
        s = self.s
        out = dict(block_kw=np.zeros(self.n), reason="", active=[], beyond=[], reach={}, kappa=0., cost=0.,
                   window_empty=False, shape=None, q_planned={}, eta={}, parts=None, B0=None, B=None)
        if self.tr <= 0:
            out["reason"] = "no_hours_left"; return out
        t = s.alpha/(1-s.alpha)
        R = list(rivals_sorted[:s.cert_rivals])
        out["active"] = R[:s.max_active]                  # logged for progress and b/J diagnostics
        B_rest = sum(self.cert_coef(j, t) for j in rivals_sorted[s.cert_rivals:])
        data = []
        for j in R:
            muL, ML, sl = self.rows(self.L, self.plan_hours)
            muj, Mj, _ = self.rows(j, self.plan_hours)
            d = self.wh(muL-muj).ravel()
            S = self.sigma2*np.eye(d.size)
            for k in (self.L, j):
                Phi = self.whp(self.cands[k].design(self.Pf[sl], self.Qf[sl]))
                S += .5*Phi@self.covs[k]@Phi.T
            cf = cho_factor(S)
            Sd = cho_solve(cf, d)
            data.append((self.cert_coef(j, t), ML-Mj, cf, Sd, float(d@Sd)))

        def bound(qs):
            return B_rest+sum(c*math.exp(-q/8.) for (c, *_), q in zip(data, qs))
        B0 = bound([x[4] for x in data])
        V, delta = s.cert_value_cny, s.delta
        best = (V*max(B0-delta, 0.), None, 0., None, 0., B0)
        for name, shape in self.shapes(R[:s.max_active]):
            s_kw = shape*self.sig_kw
            kmax, parts, unit = self.window(s_kw)
            if not np.isfinite(kmax) or kmax <= 0 or not np.isfinite(unit):
                continue
            U = np.repeat((s_kw/self.base_kw)[None, :], SPH*self.tr, axis=0)
            quad = []
            for c, D, cf, Sd, a in data:
                g = self.wh(np.einsum("tim,tm->ti", D, U)).ravel()
                quad.append((a, float(g@Sd), float(g@cho_solve(cf, g))))
            for k in np.linspace(0., kmax, s.kappa_grid+1)[1:]:
                Bk = bound([a+2*b*k+cc*k*k for a, b, cc in quad])
                obj = self.level_cost(k, unit)+V*max(Bk-delta, 0.)
                if obj < best[0]-1e-12:
                    best = (obj, name, float(k), s_kw, self.level_cost(k, unit), Bk)
        _, name, k, s_kw, cost, Bk = best
        out.update(B0=float(B0), B=float(Bk))
        if name is None:
            out["reason"] = "certified_without_block" if B0 <= delta else "no_gain"
            return out
        out.update(reason="certified" if Bk <= delta else "best_effort", shape=name, kappa=k, cost=cost,
                   window_empty=bool(Bk > delta), block_kw=k*s_kw)
        return out
\end{lstlisting}

\subsection{Strong greedy}

\begin{lstlisting}[firstnumber=380,title={\code{planner.py}, \code{Gate.plan\_bayes\_greedy}, lines 380--469}]
    def plan_bayes_greedy(self, rivals_sorted):
        """One-step greedy on the Bhattacharyya bound of the Bayes misidentification risk.

        U(u) = sum over pairs (c, c') of operational classes of
               sqrt(W_c W_c') exp(-q_cc'(u)/8),   W_c = sum of the posterior weights in class c,
        q_cc' evaluated on the classes' representatives: the Bhattacharyya bound on the
        probability that the MAP class after the next hour is wrong, with equivalent
        candidates treated as one hypothesis. The block minimises c(u) + V U(u) over the dictionary.
        For a fixed menu of blocks with frozen sensitivities, U(0) - U(A) is monotone
        submodular in the set A of executed blocks, so greedy schedules are
        near-optimal (PLANNER_REVISION.md section 10).
        """
        s = self.s
        out = dict(block_kw=np.zeros(self.n), reason="", active=list(rivals_sorted[:s.max_active]), beyond=[],
                   reach={}, kappa=0., cost=0., window_empty=False, shape=None, q_planned={}, eta={}, parts=None,
                   B0=None, B=None)
        if self.tr <= 0:
            out["reason"] = "no_hours_left"; return out
        # Operational classes (manuscript (19)) carry the prior mass: equivalent candidates are one
        # hypothesis. The leader's class is the gate's equivalence set (the stopping rule's own);
        # the other candidates with E < 25 are grouped around representatives on the same full-day
        # predicted voltages, at most bayes_members classes (2026-11-22 fix: members chosen by
        # individual weight were all twins of the leader whenever it had many equivalents).
        pool = [int(k) for k in np.argsort(self.E) if self.E[k] < 25.]
        if self.L not in pool:
            pool = [self.L]+pool
        volt = {}

        def v_of(k):
            if k not in volt:
                volt[k] = scores.voltages(self.cands[k], self.thetas[k], self.Pf, self.Qf, self.V0f,
                                          s.use_loss_term)
            return volt[k]
        members, mass = [self.L], [0.]
        for k in pool:
            if k == self.L or self.eq[k]:
                mass[0] += float(self.w[k]); continue
            vk = v_of(k)
            hit = next((i for i in range(1, len(members))
                        if float(np.max(np.abs(vk-v_of(members[i])))) <= s.equivalence_pu), None)
            if hit is not None:
                mass[hit] += float(self.w[k])
            elif len(members) < s.bayes_members:
                members.append(k); mass.append(float(self.w[k]))
        w = np.asarray(mass); w = w/w.sum()
        cls = list(range(len(members)))
        hours = self.plan_hours[:1]
        data = []
        for a in range(len(members)):
            for b in range(a):
                ha, hb_ = members[a], members[b]
                mua, Ma, sl = self.rows(ha, hours)
                mub, Mb, _ = self.rows(hb_, hours)
                d = self.wh(mua-mub).ravel()
                S = self.sigma2*np.eye(d.size)
                for k in (ha, hb_):
                    Phi = self.whp(self.cands[k].design(self.Pf[sl], self.Qf[sl]))
                    S += .5*Phi@self.covs[k]@Phi.T
                cf = cho_factor(S)
                Sd = cho_solve(cf, d)
                data.append((math.sqrt(w[a]*w[b]), Ma-Mb, cf, Sd, float(d@Sd)))
        if not data:
            out["reason"] = "one_class"; return out

        def risk(qs):
            return sum(om*math.exp(-q/8.) for (om, *_), q in zip(data, qs))
        U0 = risk([x[4] for x in data])
        V = s.cert_value_cny
        best = (V*U0, None, 0., None, 0., U0)
        for name, shape in self.shapes(rivals_sorted[:s.max_active]):
            s_kw = shape*self.sig_kw
            kmax, parts, unit = self.window(s_kw)
            if not np.isfinite(kmax) or kmax <= 0 or not np.isfinite(unit):
                continue
            U = np.repeat((s_kw/self.base_kw)[None, :], SPH, axis=0)
            quad = []
            for om, D, cf, Sd, a in data:
                g = self.wh(np.einsum("tim,tm->ti", D, U)).ravel()
                quad.append((a, float(g@Sd), float(g@cho_solve(cf, g))))
            for k in np.linspace(0., kmax, s.kappa_grid+1)[1:]:
                Uk = risk([a+2*b*k+cc*k*k for a, b, cc in quad])
                obj = self.level_cost(k, unit)+V*Uk
                if obj < best[0]-1e-12:
                    best = (obj, name, float(k), s_kw, self.level_cost(k, unit), Uk)
        _, name, k, s_kw, cost, Uk = best
        out.update(B0=float(U0), B=float(Uk), classes=len(set(cls)), members=len(members))
        if name is None:
            out["reason"] = "no_gain"; return out
        out.update(reason="bayes_greedy", shape=name, kappa=k, cost=cost, block_kw=k*s_kw)
        return out
\end{lstlisting}

\subsection{Broad check}

\begin{lstlisting}[firstnumber=594,title={\code{planner.py}, \code{check\_block}, lines 594--609}]
def check_block(settings, res, hour, soc, rng):
    """Broad zero-sum check (24): random signs at rho x headroom, buys = sells."""
    up = res["battery_kw"]+(res["ev_kw"] if sc.ev_available(settings, hour) else 0.)
    dn = np.minimum(res["battery_kw"], np.maximum((soc-settings.soc_bounds[0])*res["battery_kwh"], 0.))
    signs = rng.choice([-1., 1.], size=len(up))
    if signs.min() == signs.max():
        signs[rng.integers(len(signs))] *= -1
    u = settings.check_rho*np.where(signs > 0, up, dn)*signs
    pos, neg = u[u > 0].sum(), -u[u < 0].sum()
    if pos <= 0 or neg <= 0:
        return np.zeros_like(u)
    if pos > neg:
        u[u > 0] *= neg/pos
    else:
        u[u < 0] *= pos/neg
    return u
\end{lstlisting}

\end{document}
